\documentclass[11pt]{article}
\IfFileExists{lmodern.sty}{\usepackage{lmodern}}{}
\usepackage[T1]{fontenc}
\usepackage{amsmath,amssymb,amsthm,mathtools,geometry,hyperref,enumitem}
\geometry{margin=1.05in}
\hypersetup{pdftitle={A model of tangled type theory from overlapping generic litters},
 pdfsubject={An exposition of a consistency proof for NF},
 pdfauthor={},colorlinks=true,
 linkcolor=blue,citecolor=blue,urlcolor=blue}
\newtheorem{theorem}{Theorem}[section]
\newtheorem{lemma}[theorem]{Lemma}
\newtheorem{proposition}[theorem]{Proposition}
\newtheorem{corollary}[theorem]{Corollary}
\theoremstyle{definition}
\newtheorem{definition}[theorem]{Definition}
\newtheorem{convention}[theorem]{Convention}
\newtheorem{example}[theorem]{Example}
\theoremstyle{remark}
\newtheorem{remark}[theorem]{Remark}
\newtheorem*{idea}{Why this}
\newcommand{\NF}{\mathrm{NF}}
\newcommand{\TTT}{\mathrm{TTT}}
\newcommand{\TST}{\mathrm{TST}}
\newcommand{\Amb}{\mathrm{Amb}}
\newcommand{\Aut}{\operatorname{Aut}}
\newcommand{\Fix}{\operatorname{Fix}}
\newcommand{\cl}{\operatorname{cl}}
\newcommand{\pos}{\operatorname{pos}}
\newcommand{\ext}{\operatorname{ext}}
\newcommand{\dep}{\operatorname{depth}}
\newcommand{\Con}{\operatorname{Con}}
\newcommand{\Descr}{\operatorname{Descr}}
\newcommand{\calA}{\mathcal A}
\newcommand{\calK}{\mathcal K}
\newcommand{\calM}{\mathcal M}
\newcommand{\up}{\!\uparrow\!}
\title{A model of tangled type theory\\from overlapping generic litters\\[4pt]
\large An expository proof of the consistency of New Foundations}
\author{}
\date{29 October 2026}
\begin{document}
\maketitle

\begin{abstract}
We construct, inside ZFC+GCH, a model of Holmes's tangled type theory
\(\TTT_\omega\). By Holmes's transfer theorem this gives
\(\Con(\mathrm{ZFC})\Rightarrow\Con(\NF)\). The construction follows the
Holmes--Wilshaw strategy (a symmetric model whose types are built from
``litters'' of atoms, with extensionality enforced by a parity device) but
changes its geometry: litters are overlapping generic subsets of the atoms
rather than a partition, every parent description owns a large fibre of
litters rather than a single one, and the allowable permutations are
simply the automorphisms of a homogeneous structure with finitely many
parent coordinates at each vertex. With these changes the near-litter machinery, strong
supports and the hard form of the Freedom of Action theorem are not
needed; freedom of action is a back-and-forth argument. The local extension argument, the recursive type construction, and the
passage to full comprehension are given separately.
\end{abstract}

\tableofcontents

\section{Introduction}

\subsection{New Foundations}
Quine's set theory \(\NF\) \cite{Quine} has one sort of object and two
axioms: extensionality, and the comprehension scheme
\[
 \exists y\,\forall x\,(x\in y\leftrightarrow\varphi(x))
\]
for every \emph{stratified} formula \(\varphi\), i.e.\ every formula to
whose variables one can assign natural-number types so that \(x\in y\)
always has \(\mathrm{type}(y)=\mathrm{type}(x)+1\) and \(x=y\) always has
equal types. The formula \(x=x\) is stratified, so there is a universal
set \(V\) with \(V\in V\); the formula \(x\notin x\) is not, so Russell's
paradox is blocked. Specker showed that \(\NF\) refutes the axiom of
choice \cite{Specker53}, and the theory has no models in which membership
is well-founded. Whether \(\NF\) is consistent relative to ordinary set
theory was open from 1937 until the Holmes--Wilshaw proof
\cite{HW}, whose formal verification in Lean was completed in 2024.

\subsection{Why it is hard, and the route taken}
Specker also showed \cite{Specker62} that \(\NF\) is consistent if and only
if the simple theory of types \(\TST\) is consistent with the
\emph{ambiguity scheme}: every sentence is equivalent to the sentence
obtained by raising all its types by one. Holmes \cite{Holmes95} turned
this into a target one can build: \(\NF\) is consistent iff the
\emph{tangled type theory} \(\TTT\) is consistent, where \(\TTT\) is
\(\TST\) with the relation \(x^\beta\in y^\alpha\) allowed for \emph{every}
pair \(\beta<\alpha\) of types, not just consecutive ones. A model of
\(\TTT\) is a sequence of types in which each type behaves like a power
set of \emph{every} lower type at once. That cannot literally be true
(Cantor), so the model must be a \emph{symmetric} one: each type consists
only of the subsets of the lower types that are fixed by enough
permutations, in the manner of Fraenkel--Mostowski models. The difficulty
is to arrange a coherent family of permutation groups, one for every
chain of types, such that every symmetric subset of every lower type is
represented exactly once at every higher type. This coherence is a central difficulty in \cite{HW}.

\subsection{What is done here}
The model below keeps the two ideas of \cite{HW} that are essential: a
higher-type object is coded by a \emph{distinguished extension} at one
lower type, with its extensions at the other types derived from it through
``litters''; and the clashes between distinguished and derived extensions
are resolved by a \emph{parity} rule. The change is in the support geometry and extension argument.
\begin{itemize}[nosep]
\item Litters are the neighbourhoods of a \emph{random bipartite graph}
 between atoms and litter labels: they overlap, and fixing finitely many
 atoms fixes no litter.
\item Each parent description owns a \emph{fibre} of \(\mu\) litters, not
 one litter, and fibres of different descriptions are disjoint.
\item Each litter's parent information is recorded by finitely many unary
 \emph{coordinates}, so that the allowable group at each type is
 the automorphism group of one homogeneous structure.
\end{itemize}
With this, ``freedom of action'' (the statement that a finite coherent
configuration can be extended to an allowable permutation) is a
back-and-forth using the raw incidence extension property, and the near-litter bookkeeping
of \cite{HW} is not needed.

\subsection{How to read this}
Section~\ref{sec:logic} contains the logic: what \(\TST\), \(\TTT\) are
and why a model of \(\TTT_\omega\) gives \(\Con(\NF)\). It is standard and
can be skimmed. Section~\ref{sec:strategy} explains the shape of the
construction and the reason for each device, with a worked miniature.
Sections~\ref{sec:raw}--\ref{sec:union} are the proof. Each definition is
preceded by a short paragraph headed \emph{Why this}, which can be skipped
on a second reading. Section~\ref{sec:status} collects the main dependencies of the proof.

\section{From \texorpdfstring{\(\NF\)}{NF} to tangled type theory}
\label{sec:logic}

\subsection{\texorpdfstring{\(\TST\)}{TST} and typical ambiguity}
\(\TST\) is a many-sorted first-order theory with sorts (``types'')
indexed by \(\omega\), a membership relation \(x^n\in y^{n+1}\) between
consecutive types, and equality within each type. Its axioms are
extensionality at every positive type and comprehension
\(\exists y^{n+1}\forall x^n(x\in y\leftrightarrow\varphi(x))\) for every
formula \(\varphi\) (with parameters of any types). If \(\varphi\) is a
sentence of \(\TST\), \(\varphi^+\) is the sentence obtained by adding
\(1\) to the type of every variable. The \emph{ambiguity scheme} \(\Amb\)
is \(\{\varphi\leftrightarrow\varphi^+:\varphi\text{ a sentence}\}\).

\begin{theorem}[Specker \cite{Specker62}]\label{thm:specker}
\(\NF\) is consistent if and only if \(\TST+\Amb\) is consistent.
\end{theorem}

We use only the direction from \(\TST+\Amb\) to \(\NF\), taking
Specker's theorem as an external input. The construction below supplies
a model of tangled type theory; Ramsey's theorem and compactness then
supply the ambiguity scheme, as follows.

\subsection{Tangled type theory}
\begin{definition}[Holmes \cite{Holmes95}]
\(\TTT_\omega\) is the many-sorted theory with sorts indexed by \(\omega\),
equality within each sort, and a membership relation \(x^\beta\in y^\alpha\)
for \emph{every} pair \(\beta<\alpha\). Its axioms are: for every axiom
\(\psi\) of \(\TST\) (an instance of extensionality or comprehension) and
every strictly increasing map \(s:\omega\to\omega\), the sentence
\(\psi^{s}\) obtained by replacing each type \(n\) by \(s(n)\).
\end{definition}

So \(\TTT_\omega\) says: along every increasing chain of types, the
chain looks like a model of \(\TST\). It does \emph{not} say anything
relating the \(\beta\)-extension of \(x^\alpha\) to the
\(\gamma\)-extension of some \(y^\beta\in x\); those are independent.

\begin{theorem}[Holmes \cite{Holmes95}]\label{thm:holmes}
If \(\TTT_\omega\) has a model, then \(\NF\) is consistent.
\end{theorem}
\begin{proof}
Let \(\calM=(\tau_n)_{n\in\omega}\) be a model of \(\TTT_\omega\).
Let \(\Sigma\) be a finite set of TST sentences, and choose \(n\geq1\)
larger than every type index occurring in them. Colour each
\(n\)-element subset \(\{t_0<\cdots<t_{n-1}\}\) of \(\omega\) by the
truth values of the sentences in \(\Sigma\), interpreting formal type
\(i\) as \(\tau_{t_i}\). This is a finite colouring.

By Ramsey's theorem there is an infinite homogeneous set
\(H=\{h_0<h_1<\cdots\}\). The chain \((\tau_{h_i})_{i\in\omega}\)
is a model of TST, since every increasing translation of a TST axiom
holds in \(\calM\). For \(\varphi\in\Sigma\), its truth value along
\(h_0,\ldots,h_{n-1}\) equals its truth value along
\(h_1,\ldots,h_n\), by homogeneity. The latter is the truth value of
\(\varphi^+\) in the chain model. Thus that model satisfies
\(\varphi\leftrightarrow\varphi^+\) for every \(\varphi\in\Sigma\).

Every finite subset of \(\TST+\Amb\) is therefore satisfiable.
Compactness gives a model of \(\TST+\Amb\), and
Theorem~\ref{thm:specker} gives \(\Con(\NF)\).
\end{proof}

\subsection{What has to be built}\label{sec:target}
By Theorem~\ref{thm:holmes}, it suffices to produce sets
\(\tau_0,\tau_1,\dots\) and relations \({\in}\subseteq\tau_\beta\times
\tau_\alpha\) for \(\beta<\alpha\) such that
\begin{enumerate}[nosep,label=(\Alph*)]
\item\label{ax:ext} (extensionality) for all \(\beta<\alpha\) and
 \(x,y\in\tau_\alpha\): if \(\{z\in\tau_\beta:z\in x\}=\{z\in\tau_\beta:
 z\in y\}\) then \(x=y\);
\item\label{ax:comp} (comprehension) for every \(\TST\) formula
 \(\varphi\), every increasing chain of types interpreting its formal
 types, and all parameters, the set \(\{x:\varphi(x)\}\) at formal type
 \(i\) is \(\{z:z\in y\}\) for some \(y\) at formal type \(i+1\).
\end{enumerate}
Writing \(\ext_\beta(x)=\{z\in\tau_\beta:z\in x\}\), \ref{ax:ext} says
that \(x\mapsto\ext_\beta(x)\) is injective on \(\tau_\alpha\) for every
\(\beta<\alpha\), and \ref{ax:comp} says that enough subsets of
\(\tau_\beta\) are of the form \(\ext_\beta(x)\).

\section{The shape of the construction}\label{sec:strategy}

\subsection{Objects as distinguished extensions}
An object of type \(\alpha\) must know its extension at every lower
type. We shall code it by \emph{one} of them: a pre-object at \(\alpha\) is
a pair \((\beta,D)\) with \(\beta<\alpha\) and \(D\subseteq\tau_\beta\),
its \emph{distinguished extension}. Its extension at another type
\(\delta<\alpha\) is produced by a fixed injective map
\(\calA_\delta\) from subsets of lower types to subsets of \(\tau_\delta\):
\(\ext_\delta(\beta,D)=\calA_\delta(D)\). Injectivity of \(\calA_\delta\)
gives \ref{ax:ext} between two pre-objects whose extensions at \(\delta\)
are both \emph{derived}. It does not handle the clash between
\((\delta,Y)\), whose extension at \(\delta\) is \(Y\) itself, and
\((\beta,D)\) with \(\calA_\delta(D)=Y\).

\subsection{The parity rule}
Call \(D\) the \emph{inverse} of \(Y\) when \(Y=\calA_\delta(D)\); by
injectivity it is unique when it exists. Follow inverses:
\(Y,\ \calA^{-1}Y,\ \calA^{-1}\calA^{-1}Y,\dots\). We shall arrange that
this chain is always finite, and call its length the \emph{depth} of
\(Y\). The rule is: \emph{keep} \((\beta,D)\) iff the depth of \(D\) is
even. Then of \((\delta,Y)\) and \((\beta,\calA^{-1}Y)\), exactly one is
kept, because their depths differ by one. And every symmetric
\(Y\subseteq\tau_\delta\) is still the \(\delta\)-extension of a kept
object: of \((\delta,Y)\) if its depth is even, of \((\beta,\calA^{-1}Y)\)
otherwise. This is the Cantor--Bernstein-like device of \cite{HW}.

\subsection{Litters, and why they overlap}
What should \(\calA_\delta(D)\) be? It must be a set of type-\(\delta\)
objects determined by \(D\), injectively, with the source type recoverable.
Following \cite{HW}, \(\calA_\delta(D)\) will be a set of \emph{litters}:
special type-\(\delta\) objects \(b_\delta\), one for each label \(b\) in
a set \(B\), whose extension at the bottom type is a set \(L_b\) of atoms.
Each litter label \(b\) has a \emph{parent}, a description of an object
\(x\) at some source type \(\beta\) together with a support of \(x\); and
\(\calA_\delta(D)\) is the set of litters whose parent decodes to an
element of \(D\) with source type \(\beta\). Injectivity then needs: every
object has at least one parent description, descriptions with different
decoded values own disjoint sets of litters, and litters with source
\(\beta\) are distinguishable from litters with source \(\beta'\).

In \cite{HW} the litters partition the atoms, and each description owns
exactly one litter. A permutation of atoms that is forced, by coherence at
lower types, to move some atoms of a litter can then only map the litter
to a \emph{near-litter} (a set differing from a litter by a small set),
and the whole proof has to track near-litters. Here instead the pair
\((A,B)\) with the incidence relation \(a\in L_b\) is a \emph{random
bipartite graph}: litters overlap generically, and fixing finitely many
atoms fixes no litter label, because \(\mu\) labels share each finite
incidence pattern. Permutations will act on \(A\) and
\(B\) jointly, preserving incidence, and nothing is ever ``nearly''
preserved. Moreover each description owns a whole fibre of \(\mu\) litters,
all with the same parent; so when a permutation must send one litter to
another ``with the same parent'' there is always room.

\subsection{Symmetry, supports, and coherence along chains}
To keep every type of size \(\mu\) (so that the recursion can continue)
the types consist of \emph{symmetric} objects: those fixed by the pointwise
stabilizer of a finite \emph{support}. A support is a finite list of atoms
and litter labels, each tagged with a \emph{path}: the chain of types along
which it is seen from the top. Why paths? In \(\TTT\) the extension at
\(\gamma\) of an object \(x\) of type \(\alpha\), and the extension at
\(\gamma\) of some \(y\in\ext_\beta(x)\), are unrelated. A permutation
must therefore act on \(\ext_\gamma(x)\) and on \(\ext_\gamma(y)\) by
\emph{independent} permutations of type \(\gamma\): the first ``along
\(\gamma<\alpha\)'', the second ``along \(\gamma<\beta<\alpha\)''. So an
allowable permutation at type \(\alpha\) is really a family of raw
permutations of \(A\sqcup B\), one for each path \(\{-1,\dots,\alpha\}\),
with \(-1\) the bottom level of atoms. We realise this by giving the raw
structure one copy per path, with incidence only within copies, and letting
the group be the automorphism group of the whole thing.

The parent of a litter must be respected by these permutations: if
\(b\) has parent \((\chi,S)\) decoding to \(x\), then \(\pi(b)\) must have
a parent decoding to \(\pi(x)\). We enforce this by \emph{unary coordinate
functions} from the litter (on the path along which its parent matters)
to the entries of \(S\) (on the corresponding raised paths), together with
a \emph{colour} naming the code \(\chi\). Automorphisms preserve colours
and functions, and that is the whole coherence requirement.

\subsection{Positions}
Two finiteness properties are needed: closures under the coordinate
functions must be finite (so supports can be taken closed and the
back-and-forth works), and inverse chains for the parity rule must be
finite. Both come from one device: an injective \emph{position} map on
\(A\sqcup B\) into \(\mu\), with every litter of a fibre placed above every
position in the closure of its parent's support.

\subsection{A miniature}\label{sec:miniature}
Let \(\tau_{-1}=A\) be the atoms. Type \(0\) objects are the pre-objects
\((-1,D)\) with \(D\subseteq A\) symmetric; there is nothing to derive and
nothing to discard, so \(\tau_0\) is the set of symmetric subsets of
\(A\), including every litter \(b_0=(-1,L_b)\). At type \(1\) the
pre-objects are \((-1,D)\) with \(D\subseteq A\) and \((0,X)\) with
\(X\subseteq\tau_0\). The former has \(\ext_0=\calA_0(D)\), the set of
litters whose parent is an atom-tuple projection landing in \(D\). The
latter has \(\ext_{-1}=\varnothing\) and \(\ext_0=X\). The only possible
clash at type \(0\) is \(X=\calA_0(D)\); then \(X\) has depth \(1\), so
\((0,X)\) is discarded and \((-1,D)\) represents \(X\). The set
\(X=\tau_0\) itself is not an image (it contains non-litters), has depth
\(0\), and is represented by \((0,\tau_0)\): the universal set of type
\(1\). A permutation at type \(1\) consists of a raw permutation on the
path \(\{-1,1\}\) (acting on \(\ext_{-1}\)) and one on the path
\(\{-1,0,1\}\) (acting on the atoms inside elements of \(\ext_0\)), tied
together only through the parents of the litters on \(\{-1,0,1\}\), whose
coordinates point to atoms on \(\{-1,1\}\).

\section{Parameters and the raw structure}\label{sec:raw}

Work in ZFC+GCH and put
\[
 \kappa=\beth_\omega=\aleph_\omega,\qquad
 \mu=\kappa^+=2^\kappa=\aleph_{\omega+1}.
\]
Then \(\mu\) is regular and \(\mu^{<\mu}=\mu\): for \(1\leq\theta<\mu\),
\(\theta\leq\kappa\) and \(\mu^\theta\leq(2^\kappa)^\kappa=\mu\). The
role of \(\kappa\): all the counting in \S\ref{sec:count} happens at finite
beth levels, and \(\kappa\) is the first cardinal above all of them.

\begin{idea}
A saturated random bipartite graph gives, for free, every extension
property we shall need: any finite (indeed any \(<\mu\)) pattern of
incidences is realised \(\mu\) times, and any partial isomorphism of size
\(<\mu\) extends to an automorphism.
\end{idea}

Let \(T_0\) be the theory of the random bipartite graph: two sorts
\(A\) (atoms) and \(B\) (litter labels), a relation \(E\subseteq A\times
B\), and the extension axioms: for disjoint finite \(X,Y\subseteq B\) there
is \(a\in A\) with \(E(a,b)\) for all \(b\in X\) and for no \(b\in Y\), and
symmetrically. \(T_0\) is complete and \(\omega\)-categorical with
quantifier elimination (it is the theory of the Fra\"iss\'e limit of finite
bipartite graphs). Since \(\mu^{<\mu}=\mu\), it has a saturated model of
cardinality \(\mu\); fix one, \(R=(A,B,E)\), and write
\(L_b=\{a\in A:E(a,b)\}\). The sets \(L_b\) are nonempty and pairwise
distinct.

\begin{lemma}\label{lem:saturation}
Let \(P\subseteq A\sqcup B\), \(|P|<\mu\), and let \(p\) be a complete
non-algebraic \(1\)-type over \(P\). Then \(\mu\) elements
realise \(p\). Every partial isomorphism of \(R\) of size \(<\mu\) extends
to an automorphism.
\end{lemma}
\begin{proof}
A non-algebraic type remains consistent after excluding finitely many
individual elements. If fewer than \(\mu\) elements realised \(p\),
adding the inequalities
\(x\neq e\) for those elements would give a consistent type over fewer
than \(\mu\) parameters that is not realised, contradicting saturation.
Quantifier elimination makes every partial isomorphism elementary.
The standard back-and-forth for saturated models of regular cardinality
then gives the second claim.
\end{proof}

A \emph{litter cell} is a complete \(1\)-type of sort \(B\) over a set of
fewer than \(\mu\) atoms, i.e.\ an incidence pattern towards those atoms;
every pattern is consistent. There are \(\mu\) litter cells. A set
\(X\subseteq B\) is \emph{rich} if it meets every litter cell in \(\mu\)
points.

\begin{idea}
Reservoirs separate litters by the source and target type of their parent,
so that a set of litters reveals the type of the objects it was derived
from. The reserve \(R_*\) supplies litters with no parent at all.
\end{idea}

Partition \(B\) into rich sets \(R_{\beta\gamma}\), for \(\beta\in\{-1\}
\cup\omega\), \(\gamma\in\omega\), \(\beta\neq\gamma\), and a rich reserve
\(R_*\): enumerate the triples (cell, reservoir name, \(\nu<\mu\)) in order
type \(\mu\), at each step choose a not yet assigned label in the cell
(possible by Lemma~\ref{lem:saturation}, fewer than \(\mu\) having been
assigned), and put the remaining labels into \(R_*\). Reservoirs are
bookkeeping sets, not part of any structure below.

Finally fix an injection \(\pos:A\sqcup B\to\mu\), the \emph{position}.
Positions are not preserved by anything; they only order the construction.

\section{Tokens, codes and fibres}\label{sec:codes}

\begin{definition}[Paths, tokens, tuples]
A \emph{path} is a finite set \(p\subseteq\{-1\}\cup\omega\) with
\(-1\in p\); its \emph{type} is \(\max p\). A \emph{token at type
\(\beta\)} is a pair \((t,p)\) with \(t\in A\sqcup B\) and \(\max p=
\beta\); it is an \emph{atom token} if \(t\in A\) and a \emph{litter
token} if \(t\in B\). At type \(-1\) the only path is \(\{-1\}\) and only
atom tokens are allowed. For \(\alpha>\beta\), the \emph{raising} of a
token is \((t,p)\up\alpha=(t,p\cup\{\alpha\})\). A \emph{tuple at type
\(\beta\)} is a finite sequence of tokens at type \(\beta\), repetitions
allowed; its \emph{signature} records the sorts and paths of its entries.
\end{definition}

The whole construction is a recursion on the type; its order is spelled
out in \S\ref{sec:count}. In this section we are at type \(\beta\) and
assume that the structure \(M_\beta\) and group \(G_\beta\)
(\S\ref{sec:aux}) and the type \(\tau_\beta\) with its \(G_\beta\)-action
(\S\ref{sec:types}) are available. \(G_\beta\) acts on tokens at type
\(\beta\), hence on tuples. A tuple \(S\) \emph{supports} \(x\in
\tau_\beta\) if the pointwise stabilizer \(\Fix(S)\) fixes \(x\).

\begin{idea}
A parent must be something an automorphism can be required to preserve.
``The object \(x\)'' is not: there are \(\mu\) objects and we cannot have
\(\mu\) colours in a language whose finite diagrams we need to count. So we
record a parent as a \emph{code} (one of fewer than \(\kappa\)
equivariant functions, the colour) applied to a \emph{support tuple} (the
coordinates). Preserving colour and coordinates then forces the decoded
value to be transported correctly, by equivariance of the code.
\end{idea}

\begin{definition}[Codes]
Let \(\beta\geq0\). A \emph{code at type \(\beta\)} of signature
\(\sigma\) is a \(G_\beta\)-equivariant function \(\chi\) from \emph{all}
tuples of signature \(\sigma\) to \(\tau_\beta\) such that \(\chi(S)\) is
supported by \(S\) for every \(S\).
\end{definition}

\begin{lemma}[Totality]\label{lem:total}
Let \(\beta\geq0\), \(x\in\tau_\beta\), and let \(S_0\) support \(x\).
Define \(\chi(gS_0)=gx\) for \(g\in G_\beta\), and \(\chi(T)=
\varnothing_\beta\) (the empty object of type \(\beta\), see
\S\ref{sec:types}) for \(T\) outside the orbit of \(S_0\). Then \(\chi\)
is a code with \(\chi(S_0)=x\).
\end{lemma}
\begin{proof}
If \(gS_0=g'S_0\) then \(g^{-1}g'\) fixes \(S_0\), hence \(x\), so
\(gx=g'x\). Equivariance is immediate; \(gx\) is supported by \(gS_0\);
\(\varnothing_\beta\) is supported by everything.
\end{proof}

At \(\beta=-1\) we have \(\tau_{-1}=A\) and \(G_{-1}=\Aut(R)\). An atom
fixed by the pointwise stabilizer of a finite atom tuple \(S\) is an entry
of \(S\), by Lemma~\ref{lem:saturation} (move any other atom to one with
the same incidences and extend). The \emph{raw codes} are therefore the
coordinate projections \(\chi_i(S)=S_i\) on atom tuples, one for each
signature and position; there are countably many, none nullary, and all
have atom-only signatures. The last fact is used once, in
Lemma~\ref{lem:lower}.

\begin{definition}[Code families, descriptors]
For \(\beta\geq0\) choose, for each signature, representatives of the
\(G_\beta\)-orbits of tuples of that signature; for each representative
\(S_0\) and each \(x\in\tau_\beta\) supported by \(S_0\), put the code of
Lemma~\ref{lem:total} into \(C_\beta\). Put \(C_{-1}=\{\text{raw codes}\}
\). A \emph{descriptor at \(\beta\)} is a pair \((\chi,S)\) with \(\chi\in
C_\beta\) and \(S\) of the signature of \(\chi\); its \emph{decoded value}
is \(\chi(S)\). Write \(\Descr_\beta\) for the set of descriptors and
\(c_\beta=|C_\beta|\).
\end{definition}

\(C_\beta\) \emph{covers every supported pair}: if \(S=gS_0\) supports
\(x\) then \(g^{-1}x\) is supported by \(S_0\) and its code has
\(\chi(S)=x\). In particular every object has a descriptor. Different
descriptors may decode to the same object. Section~\ref{sec:count} proves
\(c_\beta\leq\beth_{\beta+1}<\kappa\); there are \(\mu\) tuples, so
\(|\Descr_\beta|=\mu\).

\begin{idea}
Positions of fibre litters are placed above the whole parent closure of
the support. One choice gives both finiteness properties we need: parent
edges strictly decrease position, and so do the inverse steps of the
parity rule.
\end{idea}

\begin{definition}[Descriptor positions and fibres]\label{def:fibres}
Choose an injection \(q_\beta:\Descr_\beta\to\mu\) with
\[
 q_\beta(\chi,S)>\pos(t)\quad\text{for every token }t\in\cl_\beta(S),
\]
where \(\cl_\beta\) is the parent closure in \(M_\beta\)
(\S\ref{sec:aux}; \(\cl_{-1}(S)=S\)). This is a recursion of length
\(\mu\): \(\cl_\beta(S)\) is finite, so each step excludes a bounded
initial segment of \(\mu\) and fewer than \(\mu\) used values. For every
\(\gamma\in\omega\), \(\gamma\neq\beta\), and every descriptor allocate a
\emph{fibre}
\[
 F_{\beta\gamma}(\chi,S)\subseteq R_{\beta\gamma}\cap
 \{b:\pos(b)>q_\beta(\chi,S)\},
\]
rich, the fibres of distinct descriptors being disjoint.
\end{definition}

The allocation is a recursion of length \(\mu\) through all triples
(descriptor, litter cell, \(\nu<\mu\)): at each step the cell meets
\(R_{\beta\gamma}\) in \(\mu\) labels, fewer than \(\mu\) of which are
used or have position \(\leq q_\beta(\chi,S)\). Fibres at different
source types lie in different reservoirs, so all fibres, over all types,
are pairwise disjoint, and each label lies in at most one fibre.

\begin{lemma}[Positions decrease across parent edges]\label{lem:posdec}
If \(b\in F_{\beta\gamma}(\chi,S)\) then \(\pos(b)>\pos(t)\) for every
\(t\in\cl_\beta(S)\), in particular for every entry of \(S\).
\end{lemma}

\section{The path structures and freedom of action}\label{sec:aux}

Fix \(\alpha\geq0\) and assume the fibres with source \(\gamma<\alpha\)
have been allocated. Write a path \(p\) of type \(\alpha\) with
\(|p|\geq2\) as
\[
 p=\{-1,\delta\}\cup P,\qquad\delta=\text{second smallest element of }p,
 \qquad P=p\setminus\{-1,\delta\},
\]
so \(P=\varnothing\) iff \(|p|=2\), and \(\alpha\in P\) otherwise.

\begin{idea}
A litter token \((b,p)\) stands for the litter \(b_\delta\in\tau_\delta\)
seen from type \(\alpha\) along the chain \(P\): it is an element of
\(\ext_\delta\) of some object at type \(\min P\), which is in turn an
element of an extension of an object at the next type of \(P\), and so on
up to \(\alpha\). The parent of \(b\), with source type \(\gamma\), matters
for the object at type \(\min P\) only if that object can have
distinguished extension at \(\gamma\), i.e.\ if \(\gamma<\min P\). On such
a path we record the parent; on other paths the token is \emph{flexible}.
The parent's support entries, which are tokens at type \(\gamma\), are
seen from \(\alpha\) along the same chain \(P\): hence the raised paths
\(u\cup P\).
\end{idea}

A fibre index \((\gamma,\delta')\) is \emph{active on \(p\)} iff
\(\delta'=\delta\), \(P\neq\varnothing\), \(\gamma\neq\delta\) and
\(\gamma<\min P\). Then \(\gamma<\alpha\).

\begin{definition}[The structure \(M_\alpha\)]
\(M_\alpha\) has, for every path \(p\) of type \(\alpha\), two sorts
\(A_p=A\times\{p\}\) and \(B_p=B\times\{p\}\) with the incidence relation
\(E_p\) copied from \(R\). Incidence relates only vertices on the same
path; the parent coordinates defined next may join different paths.
If \(b\in F_{\gamma\delta}(\chi,S)\) and \((\gamma,\delta)\) is active on
\(p\), the vertex \(b_p=(b,p)\) carries the \emph{colour}
\((\gamma,\delta,\chi)\) and the unary \emph{coordinate functions}
\[
 f_i(b_p)=(t_i,\,u_i\cup P)\qquad(i<k),\qquad\text{where }S=((t_i,u_i))_{i<k}.
\]
All other litter vertices are \emph{flexible} and carry no colour and no
functions; atom vertices carry no functions. Put \(G_\alpha=\Aut(M_\alpha)
\), with the topology of pointwise convergence; the stabilizers
\(\Fix(S)\) of finite sets are a base of open subgroups. At \(\alpha=-1\)
put \(M_{-1}=R\).
\end{definition}

Since \(\max u_i=\gamma<\min P\), the set \(u_i\cup P\) is a path of type
\(\alpha\); and \(f_i(b_p)\) never lies in the sort of \(b_p\), because
\(u_i\cup P=p\) would force \(\gamma=\delta\). Each label lies in at most
one fibre, so each vertex has at most one colour. \(M_\alpha\) mentions
only codes of source below \(\alpha\).

A set of vertices is \emph{closed} if it contains \(f_i(v)\) whenever
this is defined for one of its vertices \(v\); \(\cl(X)\) is the least
closed superset of \(X\). By Lemma~\ref{lem:posdec} every parent edge
\(v\mapsto f_i(v)\) strictly decreases the position of the underlying
label; every vertex has finitely many outgoing edges; so by K\"onig's lemma
\emph{the closure of a finite set is finite}. Incidence is not part of
closure.

\begin{idea}
Freedom of action says: any finite coherent configuration can be realised
and any two realisations are conjugate. To state it we must say exactly
which finite configurations are coherent. Nothing hidden should be
required beyond the visible colour, coordinate and acyclicity conditions;
this is what makes the back-and-forth go through.
\end{idea}

\begin{definition}[The class \(\calK_\alpha\)]
Let \(\calK_\alpha\) be the class of structures \(\mathfrak A\) in the
language of \(M_\alpha\) such that
\begin{enumerate}[nosep,label=(K\arabic*)]
\item incidence between \(A_p\) and \(B_p\) is arbitrary;
\item a litter vertex on path \(p\) carries at most one colour, and a
 colour \((\gamma,\delta,\chi)\) occurs only on paths on which
 \((\gamma,\delta)\) is active, with \(\chi\in C_\gamma\);
\item a vertex of colour \((\gamma,\delta,\chi)\), where \(\chi\) has
 signature \(((s_i,u_i))_{i<k}\), carries exactly the functions
 \(f_0,\dots,f_{k-1}\), with \(f_i(v)\) a vertex of sort \(s_i\) on path
 \(u_i\cup P\); uncoloured litter vertices and atom vertices carry no
 functions (so \(\mathfrak A\) is closed);
\item the parent graph \(v\mapsto f_i(v)\) is acyclic.
\end{enumerate}
\end{definition}

\begin{lemma}[Closed age and freedom of action]\label{lem:foa}
The finite closed substructures of \(M_\alpha\) are, up to isomorphism,
exactly the finite members of \(\calK_\alpha\). More precisely: if
\(D\subseteq M_\alpha\) is closed with \(|D|<\mu\), and \(\mathfrak B\in
\calK_\alpha\) contains \(D\) as a closed substructure with only finitely
many vertices outside \(D\), then \(\mathfrak B\) embeds into \(M_\alpha\)
over \(D\). Every isomorphism between closed substructures of \(M_\alpha\)
of size less than \(\mu\) extends to an automorphism.
\end{lemma}
\begin{proof}
Every finite closed substructure satisfies (K1)--(K3) by construction.
Condition (K4) follows from Lemma~\ref{lem:posdec}.

For the extension claim, place the new atoms of \(\mathfrak B\) one at a
time: each needs an incidence pattern towards the litters already present
in its sort, a type over fewer than \(\mu\) parameters, realised by
\(\mu\) atoms (Lemma~\ref{lem:saturation}); choose one not yet used in
that sort. Then place the new litters in an order in which every vertex
comes after its parents, which exists by (K4). A flexible litter needs an
incidence pattern towards all atoms now present in its sort; choose a
fresh label of \(R_*\) in that cell, which exists by richness. A litter
\(v\) of colour \((\gamma,\delta,\chi)\) has all its parents placed, say at
the tuple \(S'\) (after deleting \(P\) from their paths); choose a fresh
label in \(F_{\gamma\delta}(\chi,S')\) in the required cell, which exists
because a fibre is allocated to \emph{every} tuple of the signature of
\(\chi\), and is rich. The chosen vertex has exactly the colour and
coordinates prescribed by \(\mathfrak B\), no other functions, and the
prescribed incidences. There is no requirement to use distinct underlying
labels on different paths. This gives an embedding of \(\mathfrak B\) over \(D\).

For the last claim run a back-and-forth of length \(\mu\) through all
vertices. Adding a vertex \(v\) to a closed domain \(D\) of size \(<\mu\)
adds the finite set \(\cl(v)\); no old vertex acquires a new parent, as
\(D\) is closed; and \(D\cup\cl(v)\), transported along the current
partial isomorphism, is a member of \(\calK_\alpha\) containing the image
of \(D\) as a closed substructure with finitely many new vertices, so it
can be realised over the image. Regularity of \(\mu\) keeps all
intermediate maps of size \(<\mu\).
\end{proof}

For \(\alpha=-1\), read \(M_{-1}=R\), \(\cl(S)=S\), \(\calK_{-1}\) as the
class of bipartite graphs; Lemma~\ref{lem:foa} is then
Lemma~\ref{lem:saturation}.

\begin{lemma}[Fixed tokens]\label{lem:fixed}
For a finite set \(S\) of vertices of \(M_\alpha\) (\(\alpha\geq-1\)),
the vertices fixed by every element of \(\Fix(S)\) are exactly those of
\(\cl(S)\). In particular \(\Fix(S)=\Fix(\cl S)\), and every object with a
finite support has a finite closed support.
\end{lemma}
\begin{proof}
Automorphisms preserve the functions, so \(\cl(S)\) is fixed. Let \(v\notin
\cl(S)\) and \(D=\cl(S)\cup\cl(v)\). Form the free amalgam of \(D\) with a
disjoint isomorphic copy \(D'\) over \(\cl(S)\): identify the two copies of
\(\cl(S)\) and add no incidences. It satisfies (K1)--(K3), and (K4) because
a parent path that enters \(\cl(S)\) cannot leave it, so a cycle would lie
within one copy. By Lemma~\ref{lem:foa} it embeds into \(M_\alpha\) over
\(D\), giving \(v'\neq v\) and an isomorphism \(D\to\cl(S)\cup\cl(v')\)
fixing \(\cl(S)\) and sending \(v\) to \(v'\), which extends to
\(\pi\in\Fix(S)\) with \(\pi(v)\neq v\).
\end{proof}

\paragraph{Derivatives.}
Let \(\pi\in G_\alpha\) and \(-1\leq\beta<\alpha\). The sorts of
\(M_\alpha\) indexed by paths \(u\cup\{\alpha\}\) with \(\max u=\beta\) are
permuted by \(\pi\); deleting \(\alpha\) identifies their union with the
vertex set of \(M_\beta\). Call the induced permutation \(\pi_\beta\), the
\emph{derivative of \(\pi\) at \(\beta\)}.

\begin{lemma}[Derivatives are allowable]\label{lem:derivative}
\(\pi_\beta\in G_\beta\), and \((\pi\sigma)_\beta=\pi_\beta\sigma_\beta\).
\end{lemma}
\begin{proof}
For \(\beta=-1\), the derivative is the incidence-preserving action on
the raw copy at \(\{-1,\alpha\}\). For \(\beta\geq0\), incidence
agrees on corresponding sorts. For a path \(u=\{-1,\delta\}\cup
P'\) with \(|u|\geq3\), the active indices on \(u\) and on \(u\cup\{
\alpha\}\) coincide, since \(\min(P'\cup\{\alpha\})=\min P'\); so colours
and coordinate functions of litters on \(u\) in \(M_\beta\) and on
\(u\cup\{\alpha\}\) in \(M_\alpha\) agree after deleting \(\alpha\) from
coordinate paths. For \(u=\{-1,\beta\}\), litters carry no functions in
\(M_\beta\) but may carry colours and functions on \(\{-1,\beta,\alpha\}\)
in \(M_\alpha\). Forget these extra colours and functions before
restricting to the selected sorts. Deleting \(\alpha\) then recovers
\(M_\beta\), and the induced permutation preserves its entire structure.
\end{proof}

For a path \(p\) of type \(\alpha\), \(\pi\) permutes \(A_p\sqcup B_p\);
identifying this with \(A\sqcup B\) gives \(\pi^{p}\in\Aut(R)\). Note
\(\pi_{-1}=\pi^{\{-1,\alpha\}}\) and \((\pi_\delta)_{-1}=
\pi^{\{-1,\delta,\alpha\}}\). An automorphism of \(R\) is determined by its
action on \(A\), since distinct labels have distinct neighbourhoods.

\section{The types}\label{sec:types}

Put \(\tau_{-1}=A\) with \(G_{-1}=\Aut(R)\) acting naturally. Fix
\(\alpha\geq0\) and assume that for \(-1\leq\beta<\alpha\) the type
\(\tau_\beta\), its \(G_\beta\)-action, and the \emph{literal litters}
\(b_\beta\in\tau_\beta\) (\(b\in B\)) have been constructed, with the
properties proved in this section at the lower stages.

\begin{definition}[Alternate extensions]
For \(-1\leq\beta<\alpha\), nonempty \(X\subseteq\tau_\beta\), and
\(\delta\in\omega\) with \(\delta<\alpha\), \(\delta\neq\beta\), put
\[
 \calA_\delta(X)=\{b_\delta:\ \exists(\chi,S)\in\Descr_\beta\
 [\chi(S)\in X\ \&\ b\in F_{\beta\delta}(\chi,S)]\}\subseteq\tau_\delta,
\]
and \(\calA_\delta(\varnothing)=\varnothing\).
\end{definition}

\begin{lemma}[Injectivity]\label{lem:inj}
For nonempty \(X\subseteq\tau_\beta\), \(\calA_\delta(X)\) has cardinality
\(\mu\) and determines both \(\beta\) and \(X\).
\end{lemma}
\begin{proof}
\(\calA_\delta(X)\subseteq\{b_\delta:b\in R_{\beta\delta}\}\) determines
\(\beta\). Fibres are pairwise disjoint and of size \(\mu\), so
\(\calA_\delta(X)\) determines the set of descriptors decoding into \(X\),
and every element of \(\tau_\beta\) has a descriptor.
\end{proof}

\begin{definition}[Partial inverse, depth]
For \(0\leq\delta<\alpha\) and \(Y\subseteq\tau_\delta\) let
\(\calA^{-1}_\alpha(Y)\) be the unique \((\beta,X)\) with \(-1\leq\beta<
\alpha\), \(\beta\ne\delta\), \(X\subseteq\tau_\beta\) nonempty and \(\calA_\delta(X)=Y\), if
there is one. Iterating from \(Y\) gives a unique chain
\(Y=Y_0,Y_1,Y_2,\dots\) with \(Y_i=\calA(Y_{i+1})\); its length (the index
of the last member, which has no inverse) is the \emph{depth}
\(\dep_\alpha(Y)\), finite by Lemma~\ref{lem:depth};
\(\dep_\alpha(\varnothing)=0\). Subsets of the raw type \(\tau_{-1}\) have no inverse and have depth zero.
\end{definition}

\begin{corollary}\label{cor:literal-in-closure}
If \((\chi,S)\in\Descr_\beta\) and \(\chi(S)=c_\beta\) is a literal
litter, then \((c,\{-1,\beta\})\in\cl_\beta(S)\); hence
\(\pos(c)<q_\beta(\chi,S)<\pos(b)\) for every \(b\in F_{\beta\gamma}
(\chi,S)\).
\end{corollary}
\begin{proof}
\(c_\beta\) is supported by \(S\), and \(\pi\cdot c_\beta=(\pi_{-1}(c))_
\beta\) for \(\pi\in G_\beta\) (Lemma~\ref{lem:literal} at stage
\(\beta\)); so \(\Fix(S)\) fixes the token \((c,\{-1,\beta\})\), which
lies in \(\cl_\beta(S)\) by Lemma~\ref{lem:fixed}. Then apply
Definition~\ref{def:fibres}.
\end{proof}

\begin{lemma}[Inverse chains are finite]\label{lem:depth}
There is no infinite chain of nonempty typed sets
\(Y_0,Y_1,\dots\) with \(Y_i=\calA(Y_{i+1})\), and no such cycle.
\end{lemma}
\begin{proof}
Let \(Y_0=\calA_\delta(Y_1)\) and \(Y_1=\calA(Y_2)\). All elements of
\(Y_1\) are literal litters. Let \(b_\delta\in Y_0\) have least position
among the elements of \(Y_0\); then \(b\in F_{\beta\delta}(\chi,S)\) with
\(\chi(S)=c_\beta\in Y_1\), and \(\pos(c)<\pos(b)\) by
Corollary~\ref{cor:literal-in-closure}. So the least position strictly
decreases from \(Y_0\) to \(Y_1\), and likewise along any chain from its
second member on. Ordinals have no infinite descending sequences.
\end{proof}

\begin{definition}[Pre-objects, retained objects, \(\tau_\alpha\)]
Types are treated as separate sorts; their top-type tags are suppressed
in the notation. A \emph{pre-object at \(\alpha\)} is a pair \((\beta,D)\) with \(-1\leq
\beta<\alpha\) and \(D\subseteq\tau_\beta\); all pairs with \(D=\varnothing
\) are identified and called \(\varnothing_\alpha\). Its \emph{extensions}
are \(\ext_\beta(\beta,D)=D\); \(\ext_\delta(\beta,D)=\calA_\delta(D)\)
for \(\delta\in\omega\), \(\delta<\alpha\), \(\delta\neq\beta\); and
\(\ext_{-1}(\beta,D)=\varnothing\) if \(\beta\geq0\). The pre-object is
\emph{retained} if \(\dep_\alpha(D)\) is even. For \(\pi\in G_\alpha\) put
\(\pi\cdot(\beta,D)=(\beta,\pi_\beta[D])\). A finite tuple \(S\) at type
\(\alpha\) \emph{supports} a pre-object if \(\Fix(S)\) fixes it. Let
\(\tau_\alpha\) be the set of retained pre-objects with a finite support.
Membership: for \(z\in\tau_\beta\), \(\beta<\alpha\), \(z\in x\) iff
\(z\in\ext_\beta(x)\).
\end{definition}

\begin{idea}
The next lemma is the whole point of the colours and coordinate functions:
it says the group action commutes with \(\calA_\delta\). Without it,
\(\pi\cdot x\) would not have the right extensions and \(\tau_\alpha\)
would not be a \(G_\alpha\)-set.
\end{idea}

\begin{lemma}[Fibre transport]\label{lem:transport}
For \(\pi\in G_\alpha\), \((\chi,S)\in\Descr_\beta\) and \(\delta\) as
above,
\[
 \pi^{\{-1,\delta,\alpha\}}\bigl[F_{\beta\delta}(\chi,S)\bigr]
 =F_{\beta\delta}(\chi,\pi_\beta[S]),
\]
and consequently \(\pi_\delta[\calA_\delta(X)]=\calA_\delta(\pi_\beta[X])\)
for all \(X\subseteq\tau_\beta\). Hence \(\ext_\delta(\pi\cdot x)=
\pi_\delta[\ext_\delta(x)]\) for every \(\delta\), the map
\(\calA^{-1}_\alpha\) commutes with the action, depth is
\(G_\alpha\)-invariant, and \(G_\alpha\) acts on \(\tau_\alpha\).
\end{lemma}
\begin{proof}
On the path \(\{-1,\delta,\alpha\}\), \(P=\{\alpha\}\) and \((\beta,
\delta)\) is active, so \(b_{\{-1,\delta,\alpha\}}\) has colour \((\beta,
\delta,\chi)\) with coordinates \(S\up\alpha\) iff \(b\in F_{\beta\delta}
(\chi,S)\). As \(\pi\) preserves colours and functions, \(\pi(b)\) has
colour \((\beta,\delta,\chi)\) and coordinates \(\pi[S\up\alpha]=
\pi_\beta[S]\up\alpha\), i.e.\ \(\pi(b)\in F_{\beta\delta}(\chi,\pi_\beta
[S])\); the same for \(\pi^{-1}\) gives equality. For the second claim,
\(\chi(\pi_\beta[S])=\pi_\beta(\chi(S))\) by equivariance, so
\(\chi(\pi_\beta[S])\in\pi_\beta[X]\) iff \(\chi(S)\in X\); and
\(\pi_\delta\in G_\delta\) acts on literal litters through
\((\pi_\delta)_{-1}=\pi^{\{-1,\delta,\alpha\}}\)
(Lemma~\ref{lem:literal} at stage \(\delta\)). The remaining claims follow
with Lemma~\ref{lem:derivative}.
\end{proof}

\begin{lemma}[Literal litters]\label{lem:literal}
For \(b\in B\), \(b_\alpha:=(-1,L_b)\) is in \(\tau_\alpha\), supported by
the token \((b,\{-1,\alpha\})\), and \(\pi\cdot b_\alpha=(\pi_{-1}(b))_
\alpha\).
\end{lemma}
\begin{proof}
Subsets of \(\tau_{-1}\) have depth \(0\), since no \(\calA_\delta\) has
target \(-1\). And \(\pi\cdot(-1,L_b)=(-1,\pi_{-1}[L_b])=(-1,L_{\pi_{-1}
(b)})\) because \(\pi_{-1}\) preserves incidence; this is fixed when
\(\pi\) fixes \((b,\{-1,\alpha\})\).
\end{proof}

\begin{lemma}[Extensionality]\label{lem:ext}
If \(x,y\in\tau_\alpha\), \(\delta\in\omega\), \(\delta<\alpha\), and
\(\ext_\delta(x)=\ext_\delta(y)\neq\varnothing\), then \(x=y\). Moreover
every nonempty \(x\in\tau_\alpha\) has \(\ext_\delta(x)\neq\varnothing\)
for every proper \(\delta<\alpha\).
\end{lemma}
\begin{proof}
Write \(x=(\beta,D)\), \(y=(\beta',D')\). If \(\beta=\beta'=\delta\) then
\(D=D'\). If \(\beta,\beta'\neq\delta\) then \(\calA_\delta(D)=\calA_\delta
(D')\), and Lemma~\ref{lem:inj} gives \(\beta=\beta'\), \(D=D'\). If
\(\beta=\delta\neq\beta'\), then \(D=\calA_\delta(D')\) with \(\beta'<
\alpha\), so \(\calA^{-1}_\alpha(D)=(\beta',D')\) and \(\dep_\alpha(D)=
\dep_\alpha(D')+1\); the parities differ, so \(x,y\) are not both
retained. The last claim is Lemma~\ref{lem:inj}.
\end{proof}

Note that nothing is claimed at \(\delta=-1\): a pre-object \((\beta,D)\)
with \(\beta\geq0\) has empty extension there. Type \(-1\) is scaffolding
and is not a type of the model (\S\ref{sec:target}), exactly as in
\cite{HW}.

\begin{lemma}[Symmetric comprehension]\label{lem:symm}
Let \(-1\leq\beta<\alpha\), and let \(X\subseteq\tau_\beta\) be stabilized
setwise by \(\Fix(S)\) for a finite tuple \(S\) at type \(\alpha\) (acting
through \(\pi\mapsto\pi_\beta\)). Then there is an \(x\in\tau_\alpha\)
with \(\ext_\beta(x)=X\), supported by \(S\). If \(\beta\geq0\), this
representative is unique. In
particular \(\tau_\alpha\) contains \(\varnothing_\alpha\), every
singleton \(\{z\}_\alpha\) for \(z\in\tau_\beta\), and finite unions.
\end{lemma}
\begin{proof}
If \(\dep_\alpha(X)\) is even, take \((\beta,X)\). If it is odd, then
\(X=\calA_\beta(X')\) with \((\beta',X')=\calA^{-1}_\alpha(X)\) of even
depth; take \((\beta',X')\), whose extension at \(\beta\) is \(X\). In
either case \(\Fix(S)\) stabilizes the distinguished extension, by
Lemma~\ref{lem:transport} and injectivity, so \(S\) supports the object.
Uniqueness for \(\beta\geq0\) is Lemma~\ref{lem:ext}. For a singleton,
raise a finite support of its member to \(\alpha\). A singleton is never
an \(\calA\)-image, as fibres are infinite, so it has depth \(0\);
\(\varnothing_\alpha\) has empty support; a finite union of symmetric sets
is symmetric.
\end{proof}

\begin{remark}[Stages]\label{rem:stage}
\(\calA^{-1}_\alpha\) looks only at source types below \(\alpha\). A set
\(Y\subseteq\tau_\delta\) may acquire an inverse with source \(\beta'\geq
\alpha\) at a later stage, so its depth may change. This is harmless:
membership in \(\tau_\alpha\) is decided once, at stage \(\alpha\);
extensionality is a statement about objects of one type and uses only
\(\dep_\alpha\); and \(G_\alpha\) preserves \(\dep_\alpha\). Literal
litters, \(\varnothing_\alpha\) and singletons are retained at every stage.
\end{remark}

\begin{idea}
Comprehension for formulas whose quantifiers stay below the type of the
set being defined is now automatic, because the action preserves
membership along chains. The only subtlety is \emph{which} derivative acts
on which variable: the one along the suffix of the chain.
\end{idea}

\begin{lemma}[Predicative comprehension along chains]\label{lem:pred}
Let \(t_0<\dots<t_n\) be proper types. Let \(\varphi(z,\vec w)\) be a
\(\TST\) formula with \(z\) of formal type \(j<n\), parameters \(w_i\) of
formal types \(\leq j\), and all bound variables of formal type \(\leq j\).
Interpret formal type \(m\) as \(\tau_{t_m}\). Then for all parameters,
\(\{z\in\tau_{t_j}:\varphi(z,\vec w)\}=\ext_{t_j}(y)\) for some \(y\in
\tau_{t_{j+1}}\).
\end{lemma}
\begin{proof}
For \(\pi\in G_{t_{j+1}}\), let \(\pi\) act on a variable of formal type
\(m\leq j\) by the derivative of \(\pi\) along the suffix \(t_m<t_{m+1}<
\dots<t_{j+1}\). This preserves equality, and preserves \(u^{m}\in
v^{m+1}\): the extension at \(t_m\) of the image of \(v\) is the image of
\(\ext_{t_m}(v)\) under the derivative at \(t_m\) of the permutation acting
on \(v\) (Lemma~\ref{lem:transport}), which is the permutation acting on
\(u\). By induction on \(\varphi\), the defined set is mapped to the set
defined from the images of the parameters. Raise finite supports of the
parameters to \(t_{j+1}\) and concatenate: the pointwise stabilizer fixes
every parameter, hence stabilizes the defined set, and
Lemma~\ref{lem:symm} applies.
\end{proof}

\section{Counting, and the order of the recursion}\label{sec:count}

\begin{idea}
Everything hinges on \(|\Descr_\beta|=\mu\): with more descriptors there
would be no room for disjoint fibres of size \(\mu\) in \(B\). Descriptors
are pairs (code, tuple); there are \(\mu\) tuples, so what must be
controlled is the number of codes, which comes from the number of objects
with a given support, which comes from the number of orbits of the
stabilizer. This is the analogue of the ``counting orbits'' part of
\cite{HW}, and it is short because closures are finite and
Lemma~\ref{lem:foa} says orbits are determined by closed diagrams.
\end{idea}

Let \(o_n\) bound the number of \(G_n\)-orbits of finite tuples. By
Lemma~\ref{lem:foa}, two tuples lie in one orbit iff their closures, with
the tuple entries marked, are isomorphic; so \(o_n\) is bounded by the
number of finite marked closed diagrams in the language of \(M_n\). That
language has finitely many sorts and incidence relations, countably many
coordinate symbols, and the colours
\((\gamma,\delta,\chi)\) with \(\chi\in C_\gamma\), \(\gamma<n\). Hence
\[
 o_0\leq\aleph_0,\qquad o_n\leq\max\bigl(\aleph_0,\sup_{j<n}c_j\bigr)
 \ (n>0),\qquad c_{-1}\leq\aleph_0.
\]
For a finite tuple \(S\) at type \(0\), \(\Fix(S)\) has at most \(o_0\)
orbits on atoms, so at most \(2^{\aleph_0}\) subsets of \(A\) are
supported by \(S\); elements of \(\tau_0\) are pairs \((-1,D)\), so at
most \(\beth_1\) of them have support \(S\). Since \(C_0\) has one code per
(signature, orbit representative, supported object),
\(c_0\leq\aleph_0\cdot o_0\cdot\beth_1=\beth_1\).

Let \(n>0\) and \(S\) a finite tuple at type \(n\). By
Lemma~\ref{lem:ext}, an element of \(\tau_n\) is determined by its
extension at type \(0\), which is a union of \(\Fix(S)\)-orbits on
\(\tau_0\). Each \(z\in\tau_0\) is \(\chi(U)\) for some \(\chi\in C_0\)
and tuple \(U\) at type \(0\); if \(g\in\Fix(S)\) and \(g(U\up n)=U'\up
n\), then \(g_0(\chi(U))=\chi(U')\). So the \(\Fix(S)\)-orbit of
\(\chi(U)\) is determined by \(\chi\) and the \(G_n\)-orbit of the tuple
\(S{+}U\up n\), and there are at most \(c_0\,o_n\) orbits. Hence at most
\(2^{c_0o_n}\) elements of \(\tau_n\) have support \(S\), and
\[
 c_n\leq\aleph_0\cdot o_n\cdot2^{c_0o_n}.
\]
By induction, \(o_n\leq\beth_n\) and \(c_n\leq\beth_{n+1}\) for all
\(n\geq0\). There are \(\mu\) finite tuples, so \(|\tau_n|\leq\mu\), with
equality because of the literal litters. All code bounds stay below
\(\beth_\omega=\kappa<\mu\).

\paragraph{The recursion.}
\emph{Stage \(-1\):} \(M_{-1}=R\), \(G_{-1}=\Aut(R)\), \(\tau_{-1}=A\), the
raw codes, the positions \(q_{-1}\), the fibres \(F_{-1,\gamma}\) for all
\(\gamma\in\omega\). \emph{Stage \(n\geq0\):} form \(M_n\) from the fibres
with source below \(n\) and prove Lemmas~\ref{lem:foa}--\ref{lem:derivative}
for it; define \(\tau_n\), its action, and the literal litters, and prove
Lemmas~\ref{lem:inj}--\ref{lem:pred}; form \(C_n\), choose \(q_n\), and
allocate the fibres \(F_{n\gamma}\) for all \(\gamma\neq n\). These fibres
are inactive in every \(M_m\) with \(m\leq n\), even for \(\gamma<n\),
because a colour with source \(n\) requires a path containing an element
above \(n\). No circularity: \(M_n\) uses only codes of source below
\(n\), and \(C_n\) uses only \(M_n\) and \(\tau_n\).

\section{Singleton union and full comprehension}\label{sec:union}

\begin{idea}
Lemma~\ref{lem:pred} gives comprehension only for formulas whose
quantifiers do not go above the set being defined. \(\TST\) comprehension
has no such restriction. The standard repair, as in \cite[\S4.6]{HW}, is
to define the set of iterated singletons of the wanted elements at a high
enough type, which is predicative, and then lower it one type at a time.
The lowering step is where one needs a permutation of a \emph{lower} type
to lift to one of a \emph{higher} type fixing a given support; this is
freedom of action again, applied to a configuration made of two closed
pieces.
\end{idea}

For \(x\in\tau_\gamma\) and \(\beta>\gamma\) write \(\{x\}_\beta\in
\tau_\beta\) for the singleton of Lemma~\ref{lem:symm}. A set
\(X\subseteq\tau_\gamma\) is \emph{represented at \(\beta\)} if \(X=
\ext_\gamma(y)\) for some \(y\in\tau_\beta\); by Lemma~\ref{lem:symm}
this holds iff \(X\) has a finite \(G_\beta\)-support.

\begin{lemma}[Type lowering]\label{lem:lower}
Let \(\alpha>\beta>\gamma\geq0\) and \(X\subseteq\tau_\gamma\). If
\(\{\{x\}_\beta:x\in X\}\) is represented at \(\alpha\), then \(X\) is
represented at \(\beta\).
\end{lemma}
\begin{proof}
Let \(y\in\tau_\alpha\) with \(\ext_\beta(y)=\{\{x\}_\beta:x\in X\}\), and
let \(S\) be a finite closed support of \(y\) (Lemma~\ref{lem:fixed}). Let
\(S_{(\beta)}\) consist of the entries of \(S\) whose path has
second-greatest element \(\beta\), with \(\alpha\) deleted from their
paths; a finite tuple at type \(\beta\). We show that \(S_{(\beta)}\)
supports \(X\) in \(G_\beta\).

Let \(\pi\in G_\beta\) fix \(S_{(\beta)}\) pointwise, and let \(x\in X\).
Let \(T\) be a finite support of \(x\) at type \(\gamma\); then
\(T\up\beta\) supports \(\{x\}_\beta\), and so does \(T^*=\cl_\beta(T\up
\beta)\).

\emph{\(T^*\up\alpha\) is closed in \(M_\alpha\), with the same colours and
functions as \(T^*\).} By Lemma~\ref{lem:derivative} the only vertices
that gain functions on raising are litters on the path \(\{-1,\beta\}\).
None is in \(T^*\): the vertices of \(T\up\beta\) have paths of length
\(\geq3\), and a coordinate of a vertex of \(T^*\) lies on a path
\(u\cup P'\) with \(\max u=\gamma'<\min P'\), which is \(\{-1,\beta\}\)
only if \(u=\{-1\}\), i.e.\ the code has raw source; raw signatures
consist of atoms, so that coordinate is an atom.

\emph{The map \(\theta\), identity on \(S\) and \(t\up\alpha\mapsto
\pi(t)\up\alpha\) on \(T^*\up\alpha\), is an isomorphism between closed
substructures of \(M_\alpha\).} Its domain is closed, as a union of closed
sets. Every vertex of \(T^*\up\alpha\) lies on a path with second-greatest
element \(\beta\); so an \(S\)-vertex in the same sort as such a vertex is
in \(S_{(\beta)}\up\alpha\), fixed by \(\pi\). Hence the two definitions
agree on the overlap, and incidence between an \(S\)-vertex and a moved
vertex is preserved because \(\pi\) preserves it in \(M_\beta\) and fixes
the \(S\)-vertex. Colours and functions are preserved on \(S\) (identity)
and on \(T^*\up\alpha\) (as \(\pi\) is an automorphism of \(M_\beta\)); a
vertex in the overlap has its coordinates in both parts, fixed by \(\pi\).
\(\theta\) is injective: if \(\pi(t)\up\alpha=s\in S\), then \(s\) is in
the sort of a moved vertex, so \(s\in S_{(\beta)}\up\alpha\) is fixed by
\(\pi\), and \(t\up\alpha=s\). The image \(S\cup\pi[T^*]\up\alpha\) is
closed for the same reason as the domain.

By Lemma~\ref{lem:foa}, \(\theta\) extends to \(\pi^*\in G_\alpha\). Then
\(\pi^*\) fixes \(S\), hence \(y\); and \(\pi^*_\beta\) agrees with
\(\pi\) on \(T^*\), a support of \(\{x\}_\beta\), so
\(\pi^*_\beta\cdot\{x\}_\beta=\pi\cdot\{x\}_\beta=\{\pi_\gamma(x)\}_\beta\).
As \(\pi^*\) fixes \(y\), this singleton lies in \(\ext_\beta(y)\), so
\(\{\pi_\gamma(x)\}_\beta=\{x'\}_\beta\) for some \(x'\in X\), and
\(\pi_\gamma(x)=x'\in X\) by Lemma~\ref{lem:ext}. The same for
\(\pi^{-1}\) gives \(\pi_\gamma[X]=X\). So \(S_{(\beta)}\) supports \(X\),
and Lemma~\ref{lem:symm} represents \(X\) at \(\beta\).
\end{proof}

\begin{theorem}\label{thm:model}
\((\tau_n)_{n\in\omega}\), with \(z\in x\) iff \(z\in\ext_{\mathrm{type}
(z)}(x)\), is a model of \(\TTT_\omega\). Hence \(\Con(\mathrm{ZFC})
\Rightarrow\Con(\NF)\).
\end{theorem}
\begin{proof}
Extensionality along every pair of types is Lemma~\ref{lem:ext}. For
comprehension, let \(\varphi(x)\) be a \(\TST\) formula with \(x\) of
formal type \(i\), and let \(j\) exceed every formal type in \(\varphi\).
Fix an increasing chain of proper types interpreting the formal types up to
\(j+1\) (the model has \(\omega\) types, so every finite chain extends).
Let \(s_{i,j}(x)\) be the \((j-i)\)-fold singleton of \(x\), at formal type
\(j\). The set
\[
 \{z^{j}:\exists x^{i}\,[z=s_{i,j}(x)\wedge\varphi(x)]\}
\]
is defined by a formula whose variables all have formal type \(\leq j\)
(``\(w=\{v\}\)'' is \(\forall u\,(u\in w\leftrightarrow u=v)\)), so it is
represented at formal type \(j+1\) by Lemma~\ref{lem:pred}. Applying
Lemma~\ref{lem:lower} \(j-i\) times represents \(\{x^{i}:\varphi(x)\}\) at
formal type \(i+1\). This verifies \ref{ax:ext} and \ref{ax:comp} of
\S\ref{sec:target}, so Theorem~\ref{thm:holmes} applies.

The construction was carried out in ZFC+GCH, used only for
\(\mu^{<\mu}=\mu\) and the existence of a saturated model. Since
\(\Con(\mathrm{ZFC})\) implies \(\Con(\mathrm{ZFC+GCH})\) via the
constructible universe, \(\Con(\mathrm{ZFC})\Rightarrow\Con(\NF)\). No
inaccessible cardinal is used.
\end{proof}

\section{Dependencies and scope}\label{sec:status}

The simplification lies in the local permutation argument. The global
framework---multiple extensions, parity, derivative paths, and singleton
lowering---is inherited from Holmes--Wilshaw. The following four steps
connect the new support geometry to that framework.
\begin{enumerate}[nosep]
\item \emph{Closed-map extension (Lemma~\ref{lem:foa}).}
 Total codes allow every parent tuple. Rich fibres then realize the
 required incidences, so the stated finite diagrams are exactly the
 closed age of the concrete structure.
\item \emph{Equivariance (Lemma~\ref{lem:transport}).}
 Colours and parent coordinates transport descriptors exactly. The
 stage-local inverse and its depth therefore commute with the action.
\item \emph{Counting (\S\ref{sec:count}).}
 Closed diagrams determine tuple orbits. The resulting code bounds
 leave enough labels to allocate the next stage's fibres.
\item \emph{Type lowering (Lemma~\ref{lem:lower}).}
 Atom-only raw signatures prevent a lifted singleton support from
 acquiring new litter coordinates. Its action can therefore be combined
 with the identity on the higher support and extended by the first step.
\end{enumerate}

The external inputs are the standard existence and homogeneity facts
for the saturated random bipartite graph, Specker's ambiguity theorem,
Ramsey's theorem, compactness, and the constructible-universe consistency
reduction.

The ordered variant adds sort orders, structural Ramsey theory and
fixed-point arguments to study universal ordering and Boolean prime
ideals. Those additions are separate from the consistency proof given
here.

\begin{thebibliography}{9}
\bibitem{Quine} W.\,V. Quine, \emph{New Foundations for Mathematical
Logic}, Amer.\ Math.\ Monthly 44 (1937), 70--80.
\bibitem{Specker53} E. Specker, \emph{The axiom of choice in Quine's New
Foundations for Mathematical Logic}, Proc.\ Nat.\ Acad.\ Sci.\ USA 39
(1953), 972--975.
\bibitem{Specker62} E. Specker, \emph{Typical ambiguity}, in: Logic,
Methodology and Philosophy of Science (E. Nagel, P. Suppes, A. Tarski,
eds.), Stanford Univ.\ Press, 1962, 116--124.
\bibitem{Jensen} R.\,B. Jensen, \emph{On the consistency of a slight(?)
modification of Quine's NF}, Synthese 19 (1969), 250--263.
\bibitem{Holmes95} M.\,R. Holmes, \emph{The equivalence of NF-style set
theories with ``tangled'' type theories; the construction of
\(\omega\)-models of predicative NF (and more)}, J.\ Symbolic Logic 60
(1995), 178--190.
\bibitem{HW} M.\,R. Holmes and S. Wilshaw, \emph{New Foundations is
consistent}, arXiv:1503.01406; author-hosted manuscript,
\url{https://randall-holmes.github.io/Nfproof/maybedetangled2-for-referee-clean.pdf}.
\bibitem{Forster} T.\,E. Forster, \emph{Set Theory with a Universal Set},
2nd ed., Oxford Logic Guides 31, Oxford Univ.\ Press, 1995.
\end{thebibliography}
\end{document}
